\section{三个角色各自做到哪一步}

\subsection{军师：算和说分开，判断留在代码里}

军师是这次的重心。它的工作方式是把“算”和“说”切开：伤害、行动顺序、合法性全部由游戏引擎枚举，模型不参与计算，只负责把比较结果讲成玩家能理解的取舍。

玩家面对一个局面时，引擎枚举双方所有合法行动，对每一种组合跑一遍共享结算器，得到平均收益和最坏分支；双方同速时，两种出手顺序各算一次，这样搜索不会偷偷利用随机数。模型拿到的是一份算好的事实包，里面有数字、每张战术卡的适用条件和反例。它的任务只是解释。

这个分工为什么必要，一个例子就够。同一个残血目标，火花直接伤害 52，追猎 75，两个都够击倒，都耗 2 点能量。只看面板，很自然会告诉玩家“选追猎”。但这是错的：追猎的额外伤害来自“目标已灼烧”这个前提，对方吃药、换宠、防御都会让前提失效。这种判断交给模型是猜，交给枚举是算。

“该不该调工具”这件事，我最后也没让模型做。证据包里本来就带着当前局面和最近回合事件，“我还剩多少血”“对面还有几瓶药”这类问题的答案已经在包里，不值得为它多跑一次往返；真正需要查的只有证据包结构上不含的四类：指定回合的原始事件、两个具体行动的分支模拟、整局分页统计、以及包里没有的战术规则。判定需要时，代码直接执行第一步，之后才问模型要不要继续查；判定不需要时，整个工具循环跳过。这是三轮真实评测之后才改的，过程在“难点”一节。

局内主动提示同样归军师，不归陪练——这一点我改过一次。触发条件写在代码里，共四类：宠物倒下；犹豫不决（在两个以上选项之间来回至少 3 次、悬停累计 8 秒）；明显策略错误（真实枚举出来的分差大于 5，且确实造成了后果）；长停留（同一个选项上停满 10 秒且指针还在选项区）。另有第五个信号，它自己不开口，但会改写提示的措辞与优先级：血量低于 35\% 而背包里还有回复药。安静档和“本局不再提醒”永远优先，不被任何推断覆盖。

\subsection{老师：复盘只看当时能看到的}

复盘分整局和单回合两层。整局先给行动分布和回合数，再挑出造成减员或血量大幅变化的回合；单回合用回合开始时的局面做条件比较，比如“当时目标 78 血，火花不防御打 24、防御打 8，单次攻击不足以收尾”。

有一条规则是硬的：复盘一律使用回合开始时的公开状态，不使用对手实际打出的那一招。理由很直白——如果拿事后结果当依据，任何一次失败都能被解释成“你当时选错了”，这既不公平，也没法教人。这条有测试保证：改掉对手的实际行动，复盘结论不变。

练习跟着复盘走。玩家在某个回合做了选择，教练就这个选择对应的知识点出一道题（灼烧追击、防御节奏、换宠承伤、先制与速度四类）。出题、等待作答、判题、取消这四个状态由程序控制，模型不能提前报答案，也不能把题改写成别的。题目和玩家刚才的真实选择挂钩，不是随机题库。

\subsection{陪练：先只有设计，后来补上了实现}

这一块我改动最多，先说现状。

陪练的状态不是模型编出来的，是从真实对战历史派生的：连败几局、最近有没有赢、这台机器上有没有历史记录。从状态派生出语气档位，共四档，每档有字数上限和问句上限。档位不是装饰——它会作为证据写进模型的输入，档位不同，输出的长度和语气跟着变。有两条底线：本机没有任何真实记录时返回最低档并且不编造过去；某个档位需要的事实一条都拼不出来时，明说降到最短承接句，而不是硬凑一句符合档位的漂亮话。

主动开口只挂在两个事件上，每局各一次：本局第一个同伴倒下，以及整局结束。加起来一局最多两次。

教学账本也是这一段补的。已经教过的知识点，没有新的反证就不再重复讲；但军师在局内反复观察到同一个知识点上的失误（3 次独立行动都不合理）时，会把这一课标回未掌握，老师于是有机会再讲一次，而且说得出“为什么又教”。叉掉是分角色的：叉掉教学类，教学安静但军师类照常；两类都被叉掉，才整体静默。这个分层是刻意的——玩家点叉的意思通常是“这个别讲了”，不是“你闭嘴”。

情绪这件事我做得比设计文档保守。原因是我对“给游戏助手加情绪”本身有保留：项目调研过的灵宝资料里有一条我认同的判断，情绪陪伴很容易变成油腻，玩家不打开聊天，不代表需要被安慰。所以我给它加的约束比一般助手更严——允许沉默、不主动提问、不空泛安慰、不因为一次失败就弹话。落到实现里就是那四档语气加字数上限，不是一套情绪状态机。所以如果问“陪练有没有情绪”，准确答案是：有一层受约束的语气调节，没有情绪状态机。
